\section{Conclusion}

In this paper, we identify diverse and compounding failure modes during long CoT reasoning as a central barrier to robust vision-language reasoning. To address this challenge, we propose HopChain, a scalable framework for synthesizing multi-hop vision-language reasoning data for RLVR training, where earlier hops establish the instances, sets, or conditions needed for later hops, forcing repeated visual re-grounding throughout training while terminating in specific, unambiguous numerical answers suitable for verifiable rewards. By augmenting the original RLVR data with HopChain-synthesized multi-hop data, we obtain broad and generalizable gains on 20 out of 24 benchmarks on both Qwen3.5-35B-A3B and Qwen3.5-397B-A17B. Importantly, these gains arise even though the synthesized data is benchmark-agnostic rather than designed for any specific benchmark. Additional analyses further show that preserving full chained queries is important, that multi-hop data substantially strengthens long-CoT vision-language reasoning, and that the synthesized data spans a broad difficulty range while enabling trained models to correct a broad range of errors. These results establish HopChain as an effective and scalable framework for improving generalizable vision-language reasoning in VLMs.

Despite these gains, the current pipeline still depends on successful instance segmentation, so images with no detectable objects (and thus no SAM3-segmentable instances) cannot be processed and are excluded from the current synthesis workflow. A natural next step is to reduce this dependency by introducing complementary data-construction routes for images with few or no segmentable objects, while preserving the core design principle of chained visual grounding for long-CoT RLVR training.
